%%%%%%%%%%%%%%%%%%%%%%%%%%%%%%%%%%%%%%%%%%%%%%%%%%%%%%%%%%%%%%%%%%%%%%%%%%%%%%%
% Plantilla para un artículo LaTeX en español.
% Original catalogue starter based on the public Overleaf template by
% Steph Migoni: https://www.overleaf.com/latex/templates/plantilla-para-un-
% articulo-latex/ycmrsxnswcct
%%%%%%%%%%%%%%%%%%%%%%%%%%%%%%%%%%%%%%%%%%%%%%%%%%%%%%%%%%%%%%%%%%%%%%%%%%%%%%%
\documentclass[a4paper,11pt]{article}

\usepackage[spanish,es-nodecimaldot]{babel}
\usepackage[utf8]{inputenc}
\usepackage[T1]{fontenc}
\usepackage[a4paper,top=3cm,bottom=2cm,left=3cm,right=3cm]{geometry}
\usepackage{amsmath,amsthm,amssymb}
\usepackage{booktabs}
\usepackage{graphicx}
\usepackage{xcolor}
\usepackage[colorlinks=true,allcolors=blue]{hyperref}

\title{Una guía breve para comunicar resultados experimentales}
\author{Nombre de la autora\\
  \small Departamento de Física\\
  \small nombre@ejemplo.edu\\
  \small Ciudad de México}
\date{\today}

\begin{document}
\maketitle

\begin{abstract}
  Este ejemplo muestra cómo organizar un artículo corto en español. Incluye
  una pregunta de investigación, una figura, una tabla, ecuaciones,
  referencias cruzadas y una bibliografía para que pueda sustituirse por el
  contenido propio de la autora.
\end{abstract}

\section{Introducción}
La escritura científica se beneficia de una estructura explícita: primero se
plantea una pregunta, después se describe el método y finalmente se presentan
los resultados. En física, la notación matemática permite expresar esas
relaciones con precisión; por ejemplo, la propagación de una señal puede
modelarse mediante una velocidad media $v=\Delta x/\Delta t$. Una introducción
clara también sitúa el trabajo en la literatura existente~\cite{griffiths2018}.

\section{Método}
Consideremos una medición de la posición $x$ de un móvil en distintos instantes
$t$. Ajustamos una recta a los datos y estimamos la incertidumbre de la
pendiente. La ecuación de trabajo es
\begin{equation}
  x(t)=x_0+v_0t+\frac{1}{2}at^2,
  \label{eq:cinematica}
\end{equation}
 donde $x_0$ es la posición inicial, $v_0$ la velocidad inicial y $a$ la
aceleración. La ecuación~\eqref{eq:cinematica} ilustra cómo etiquetar y citar
una expresión dentro del texto.

\subsection{Representación de una figura}
La Figura~\ref{fig:esquema} es un esquema portátil, sin archivos externos.
Para usar una imagen propia, puede reemplazarse el contenido de la caja por
\verb|\includegraphics[width=0.5\textwidth]{archivo}|.
\begin{figure}[ht]
  \centering
  \setlength{\fboxsep}{12pt}
  \fbox{\rule{0pt}{3cm}\rule{0.72\textwidth}{0pt}}
  \caption{Espacio reservado para un esquema experimental o una fotografía.}
  \label{fig:esquema}
\end{figure}

\section{Resultados}
La Tabla~\ref{tab:mediciones} resume un conjunto pequeño de observaciones.
Los datos reales pueden sustituirse directamente en el entorno \verb|tabular|.
\begin{table}[ht]
  \centering
  \caption{Ejemplo de mediciones de posición.}
  \label{tab:mediciones}
  \begin{tabular}{@{}ccc@{}}
    \toprule
    $t$ (s) & $x$ (m) & $u(x)$ (m) \\
    \midrule
    0.0 & 0.02 & 0.01 \\
    0.5 & 0.39 & 0.01 \\
    1.0 & 1.01 & 0.02 \\
    1.5 & 1.86 & 0.02 \\
    \bottomrule
  \end{tabular}
\end{table}

Una comprobación útil es comparar el ajuste con la predicción teórica. Para
este ejemplo, la suma de cuadrados puede escribirse como
\begin{equation}
  S=\sum_{i=1}^{N}\left[x_i-x(t_i)\right]^2.
  \label{eq:minimos}
\end{equation}
La notación vectorial y matricial también es frecuente:
\[
  \mathbf{A}=\begin{pmatrix}a&b\\c&d\end{pmatrix},\qquad
  \det(\mathbf{A})=ad-bc.
\]

\section{Discusión}
Los resultados deben interpretarse junto con sus incertidumbres y las
limitaciones del método. Una lista de comprobación puede ayudar:
\begin{enumerate}
  \item declarar las unidades y el número de cifras significativas;
  \item explicar qué supuestos entran en el modelo;
  \item comparar el resultado con una referencia independiente.
\end{enumerate}
También es posible usar listas sin numeración:
\begin{itemize}
  \item revisar los ejes de cada figura;
  \item describir todos los símbolos de una tabla;
  \item enlazar cada afirmación importante con una fuente.
\end{itemize}

\section{Conclusión}
Una plantilla no sustituye la argumentación científica, pero hace visibles sus
componentes principales. La ecuación~\eqref{eq:minimos}, la Figura
~\ref{fig:esquema} y la Tabla~\ref{tab:mediciones} forman un esqueleto que
puede ampliarse con secciones, datos y material suplementario. Para el origen
histórico de la formulación de campo, véase Maxwell~\cite{maxwell1865}.

\begingroup
\small
\bibliographystyle{abbrv}
\bibliography{references}
\endgroup

\end{document}
